\section{Spacetime Recognition action}
\label{sec:spacetime-action}

Section~\ref{sec:thermodynamic-recognition-square} removes the formerly missing interface between thermodynamic response and spacetime field theory.  A local-equilibrium field $\Theta$ and the response jet $j^N\Lambda$ define the pulled-back carrier $\cH^N_\Theta$.  The observation and target maps define the quotient bundle
\[
E_{\RK}=\Ker\cA/(\Ker\cA\cap\Ker\Pi),
\]
the distinguished thermodynamic Recognition field $\Phi_{\mathrm{th}}$, and the induced quotient connection $\nabla^{\RK}$.  Its curvature is the leading target-faithful transport-square defect.  Thus the fibre, field, connection, and curvature used below are constructed Recognition objects rather than a scalar and gauge field renamed after their equations are written.

Let $(M,g)$ be an oriented, time-oriented Lorentzian manifold of dimension $n\ge3$.  The metric may be supplied as calibrated spacetime data or reconstructed from the normalized scalar principal symbol of target-faithful quotient propagation by Theorem~\ref{thm:principal-symbol-recognition-metric}; characteristic data alone determine its conformal class.  Assume either that $M$ has no boundary or that all variations below have compact support.  Let $P\to M$ be a compact orthogonal or unitary frame reduction of $E_{\RK}$ that preserves its fibre metric, target structure, and any declared cut sectors.  Write its compact structure group as $G_{\RK}$, its Lie algebra as $\mathfrak g_{\RK}$, and equip $\mathfrak g_{\RK}$ with an $\operatorname{Ad}$-invariant positive inner product $\langle\cdot,\cdot\rangle_{\mathfrak g}$.  Equivalently,
\[
E_{\RK}=P\times_\rho H_{\RK},
\]
where $H_{\RK}$ is the typical fibre of the target-relevant quotient or a declared Hilbert completion of it.

The quotient connection $\nabla^{\RK}$ is represented in a local frame by a principal connection $A$ on $P$ and induces the covariant derivative $D_A$ on $E_{\RK}$.  Its curvature
\[
F_A\in\Omega^2(M;\operatorname{ad}P)
\]
is the local representative of the quotient curvature in \eqref{eq:quotient-curvature}.  A spacetime Recognition field is a section
\[
\Phi\in\Gamma(E_{\RK}).
\]
The section $\Phi_{\mathrm{th}}$ in \eqref{eq:thermodynamic-recognition-field} is the response-generated configuration singled out by the thermodynamic data; allowing $\Phi$ to vary gives the dynamical completion of that same target-relevant carrier.  Let $V_{\RK}:E_{\RK}\to\R$ be a smooth $G_{\RK}$-invariant fibre potential.

\begin{definition}[Action-lifted Recognition datum]
An action-lifted Recognition datum is
\[
\mathfrak R_{\mathrm{act}}
=(M,g;\Theta,\cH^N_\Theta,\cA,\Pi;
 P,A;E_{\RK},\Phi;V_{\RK};
 \kappa,g_{\RK},\Lambda_{\mathrm{cos}}),
\]
where $\kappa>0$, $g_{\RK}>0$, and $\Lambda_{\mathrm{cos}}\in\R$.  Its minimal parity-even local action, classified by Theorem~\ref{thm:minimal-action-classification}, is
\begin{equation}
\begin{split}
S_{\mathrm{ARK}}[g,A,\Phi]
:=\int_M\Bigg[
&\frac{1}{2\kappa}(R_g-2\Lambda_{\mathrm{cos}})
-\frac{1}{4g_{\RK}^2}
 \langle F_{A\,\mu\nu},F_A^{\mu\nu}\rangle_{\mathfrak g}
\\
&-\frac12\langle D_{A\,\mu}\Phi,D_A^{\mu}\Phi\rangle_{E_{\RK}}
-V_{\RK}(\Phi)
\Bigg]\,\mathrm{vol}_g.
\end{split}
\label{eq:ark-action}
\end{equation}
\end{definition}

For $\xi\in\mathfrak g_{\RK}$, define the Recognition current by
\begin{equation}
\langle J_\Phi^\nu,\xi\rangle_{\mathfrak g}
:=\operatorname{Re}\langle D_A^\nu\Phi,\rho_*(\xi)\Phi\rangle_{E_{\RK}}.
\label{eq:recognition-current}
\end{equation}
Define
\begin{align}
T^{A}_{\mu\nu}
&:=\frac{1}{g_{\RK}^2}
\left(
\langle F_{A\,\mu\alpha},F_{A\,\nu}{}^{\alpha}\rangle_{\mathfrak g}
-\frac14g_{\mu\nu}
 \langle F_{A\,\alpha\beta},F_A^{\alpha\beta}\rangle_{\mathfrak g}
\right),
\label{eq:gauge-stress}\\
T^{\Phi}_{\mu\nu}
&:=\operatorname{Re}\langle D_{A\,\mu}\Phi,D_{A\,\nu}\Phi\rangle_{E_{\RK}}
-g_{\mu\nu}
\left(
\frac12\langle D_{A\,\alpha}\Phi,D_A^{\alpha}\Phi\rangle_{E_{\RK}}
+V_{\RK}(\Phi)
\right).
\label{eq:recognition-stress}
\end{align}

The thermodynamic Recognition square fixes the meaning of $E_{\RK}$, $\Phi_{\mathrm{th}}$, $A$, and $F_A$.  The target-faithful quotient propagation symbol fixes $g$ when its normalization is declared.  The minimal-action theorem fixes the action form within its declared locality, symmetry, parity, and derivative-order class.  The numerical couplings and nonlinear invariant potential remain additional constitutive data; they are not inferred from equilibrium identities alone.
